\begin{ex}%[2P2B2-1][Loai1]
Hai dây đồng hình trụ có cùng khối lượng và ở cùng nhiệt độ. Dây A dài gấp đôi dây B. Điện trở dây A và điện trở dây B có mối quan hệ.
\choice
{$R_A = R_B$}
{$R_A = 2R_B$}
{$R_A = \dfrac{R_B}{4}$}
{\True $R_A = 4R_B$}
\loigiai{
Ta có: $R=\dfrac{\rho \ell}{S}\Rightarrow \dfrac{R_A}{R_B}=\dfrac{\ell_A}{\ell_B}=2\ \Omega$.
}
%<MyLT>
\end{ex}
\begin{ex}%[2P2B2-1][Loai1]
Nếu cả chiều dài lẫn đường kính của một sợi dây đồng tiết diện tròn được tăng lên gấp đôi thì điện trở của dây đó sẽ
\choice
{Không đổi}
{Tăng gấp đôi}
{\True Giảm hai lần}
{Tăng gấp bốn}
\loigiai{
Ta có: $R=\dfrac{\rho \ell}{S}=\dfrac{\rho \ell}{\pi r^2},\ R'=\dfrac{2\rho \ell}{S'}=\dfrac{2\rho \ell}{\pi {(2r)}^2}=\dfrac{\rho \ell}{\pi 2r^2}=\dfrac{R}{2}$.
}
%<MyLT>
\end{ex}
\begin{ex}%[2P2B2-1][Loai1]
Hai thanh kim loại có điện trở hoàn toàn bằng nhau. Thanh A chiều dài $\ell_A$, đường kính $d_A$, thanh
B chiều dài $\ell_B = 2\ell_A$ và đường kính $d_B = 2d_A$. Điều này suy ra rằng thanh A có điện trở suất liên hệ với thanh B theo hệ thức nào dưới đây?
\choice
{$\rho_A = \dfrac{\rho_B}{4}$}
{$\rho_A = 2\rho_B$}
{\True $\rho_A = \dfrac{\rho_B}{2}$}
{$\rho_A = 4\rho_B$}
\loigiai{
Do $R_A=R_B$ nên $\dfrac{\rho_A \ell_A}{\pi r_A^2}=\dfrac{\rho_B \ell_B}{\pi r_B^2}=\dfrac{\rho_B2 \ell_A}{\pi {(2r_A)}^2}\Rightarrow \rho_A=\dfrac{\rho_B}{2}$.
}
\end{ex}

